\documentclass[11pt]{article}

\usepackage{amsmath,amssymb,amsthm}
\usepackage[margin=1in]{geometry}

\title{Network Equilibrium and Greedy Network Expansion}
\author{}
\date{}

\begin{document}

\maketitle

\section{Sets and Parameters}

Let $N$ be the set of markets and $E$ the set of all possible undirected transmission edges.
Let $E_t$ denote the network at greedy iteration $t$.

For each market $i\in N$:

\begin{itemize}
    \item $g_i$ = production quantity,
    \item $q_i$ = consumption quantity,
    \item $S_i(q)$ = supply price function,
    \item $D_i(q)$ = demand price function.
\end{itemize}

For each transmission direction $i\rightarrow j$:

\[
\lambda_{ij}\in(0,1]
\]

is the fraction of transmitted energy received.

For each edge $e=(i,j)$:

\[
c_e = \text{construction cost}.
\]

The total construction budget is

\[
B = 10\max_{e\in E}c_e.
\]

For 11 markets,

\[
|E|=\binom{11}{2}=55.
\]

With 11 initial edges,

\[
|E_0|=11,\qquad |E\setminus E_0|=44.
\]

\section{Step-Function Blocks}

For supply block $k$ at market $i$, let

\[
K^S_{ik}=\text{capacity},\qquad
p^S_{ik}=\text{price},
\]

and let

\[
x^S_{ik}\in[0,K^S_{ik}]
\]

be the quantity used from the block.

For demand blocks,

\[
K^D_{ik}=\text{capacity},\qquad
p^D_{ik}=\text{price},
\]

with

\[
x^D_{ik}\in[0,K^D_{ik}].
\]

Total production and consumption are

\[
g_i=a^S_i+\sum_k x^S_{ik},
\]

\[
q_i=a^D_i+\sum_k x^D_{ik},
\]

where $a^S_i$ and $a^D_i$ are the minimum observed quantities.

\section{Sequential Block Constraints}

Let

\[
z^S_{ik},z^D_{ik}\in\{0,1\}
\]

indicate whether a block is completed.

For supply:

\[
x^S_{ik}\ge K^S_{ik}z^S_{ik},
\]

\[
x^S_{i,k+1}\le K^S_{i,k+1}z^S_{ik}.
\]

For demand:

\[
x^D_{ik}\ge K^D_{ik}z^D_{ik},
\]

\[
x^D_{i,k+1}\le K^D_{i,k+1}z^D_{ik}.
\]

These constraints force blocks to be filled sequentially.

\section{Network Flows}

For every active edge $(i,j)\in E_t$, define

\[
f_{ij}\ge0,\qquad f_{ji}\ge0
\]

and a binary direction variable

\[
y_{ij}\in\{0,1\}.
\]

The direction constraints are

\[
f_{ij}\le My_{ij},
\]

\[
f_{ji}\le M(1-y_{ij}),
\]

where $M$ is sufficiently large.

Thus, an edge can carry flow in at most one direction.

\section{Market Balance}

For every market $i$,

\[
\boxed{
g_i+
\sum_{j:(j,i)\in E_t}
\lambda_{ji}f_{ji}
=
q_i+
\sum_{j:(i,j)\in E_t}
f_{ij}
}
\]

Production plus received imports equals consumption plus exports.

\section{Inner Network Equilibrium}

For a fixed network $E_t$, the Gurobi model solves

\[
\boxed{
W(E_t)=
\max
\sum_{i\in N}
\left[
\sum_k p^D_{ik}x^D_{ik}
-
\sum_k p^S_{ik}x^S_{ik}
\right]
}
\]

subject to the block, quantity, flow, direction, and market-balance constraints above.

Equivalently,

\[
W(E_t)
=
\max
\left[
\sum_i B_i(q_i)
-
\sum_i C_i(g_i)
\right].
\]

\section{Autarky}

For each market, the autarky quantity satisfies

\[
S_i(q_i^A)=D_i(q_i^A).
\]

Autarky surplus is

\[
\boxed{
W_A=
\sum_{i\in N}
\left[
B_i(q_i^A)-C_i(q_i^A)
\right].
}
\]

\section{Greedy Network Expansion}

Let the current network be $E_t$ and its optimized surplus be

\[
W_t=W(E_t).
\]

For each candidate edge

\[
e\in E\setminus E_t,
\]

temporarily add the edge:

\[
E_t^e=E_t\cup\{e\}.
\]

Solve the same inner equilibrium problem to obtain

\[
W_t^e=W(E_t\cup\{e\}).
\]

The marginal surplus gain is

\[
\boxed{
\Delta W_t(e)
=
W(E_t\cup\{e\})-W(E_t).
}
\]

An edge is eligible only if

\[
\boxed{
C_t+c_e\le B
}
\]

and

\[
\boxed{
\Delta W_t(e)>0.50.
}
\]

The selected edge is

\[
\boxed{
e_t^*
=
\arg\max_{e}
\Delta W_t(e)
}
\]

over all eligible candidate edges.

The network is then updated:

\[
\boxed{
E_{t+1}=E_t\cup\{e_t^*\}
}
\]

and

\[
\boxed{
C_{t+1}=C_t+c_{e_t^*}.
}
\]

The new equilibrium surplus is

\[
\boxed{
W_{t+1}=W(E_{t+1}).
}
\]

The process terminates when no affordable edge satisfies

\[
\Delta W_t(e)>0.50,
\]

or when the maximum number of iterations is reached.

\section{Combined Model}

The exact network-design problem represented by the nested formulation is

\[
\boxed{
\begin{aligned}
\max_{E'\subseteq E\setminus E_0}
\quad&
W(E_0\cup E')
\\
\text{s.t.}\quad&
\sum_{e\in E'}c_e+
\sum_{e\in E_0}c_e
\le B,
\end{aligned}
}
\]

where $W(E)$ is itself the solution of the inner network-equilibrium optimization.

The implemented method approximates this outer problem using the greedy rule

\[
\boxed{
e_t^*
=
\arg\max_e
\left[
W(E_t\cup\{e\})-W(E_t)
\right]
}
\]

subject to the budget constraint and the minimum surplus-gain condition

\[
\boxed{
W(E_t\cup\{e\})-W(E_t)>0.50.
}
\]

\end{document}
